% The two structural operations used to assemble coefficient-triangular tables.
\begin{tikzpicture}[x=1cm,y=1cm,>=Latex]
  \node[fp panel title] at (0,4.1) {(a) Shift through a monic factor};
  \node[fp active, minimum width=42mm] (old) at (2.15,2.55)
    {old parameters $\eta_{e-1},\ldots,\eta_0$};
  \node[fp known, minimum width=24mm] (factor) at (2.15,3.42)
    {$L=x^h+\cdots$};
  \node[fp label, right=1mm of factor] {leading coefficient $1$};
  \node[fp active, minimum width=42mm] (shifted) at (1.55,1.05)
    {reindexed $\alpha_{h+e-1},\ldots,\alpha_h$};
  \node[fp unread, minimum width=16mm, right=0pt of shifted]
    {rows $<h$};
  \draw[fp decode] (old.south) -- node[fp label, right]
    {$\times L$: shift by $h$} (shifted.north);
  \node[fp label, anchor=west, text=fpBlue!80!black] at (0,0.22)
    {monicity preserves every pivot slope $\lambda_i$};

  \node[fp panel title] at (7.15,4.1) {(b) Concatenate descending blocks};
  \draw[fp decode] (7.15,3.42) -- node[fp label, above]
    {decoding order} (13.55,3.42);
  \node[fp active, minimum width=22mm, anchor=west] (tail) at (7.15,2.55)
    {new tail};
  \node[fp recursive, minimum width=27mm, anchor=west] (inner) at (9.35,2.55)
    {shifted recursion};
  \node[fp second, minimum width=18mm, anchor=west] (low) at (12.05,2.55)
    {low block};
  \draw[fp dependence] (tail.south) to[bend right=12]
    node[fp label, below] {decoded side data} (inner.south);
  \draw[fp dependence] (inner.south) to[bend right=12] (low.south);
  \node[fp label] at (8.25,1.18)
    {$I_{s-1}$ decoded first};
  \node[fp label] at (10.7,1.18)
    {then $I_{s-2},\ldots$};
  \node[fp label] at (12.95,1.18)
    {finally $I_0$};
  \node[fp label, anchor=west, text=fpBlue!80!black] at (7.15,0.22)
    {higher-block formulas may be substituted into every lower block};
\end{tikzpicture}
